% Standalone source; the open editor compiles this file in place.
% Research status: frozen prediction-input control completed, 25 October 2026. No novelty claim.
\documentclass[11pt,a4paper]{article}
\usepackage[T1]{fontenc}
\usepackage[utf8]{inputenc}
\usepackage[english]{babel}
\usepackage[margin=1in]{geometry}
\usepackage{amsmath,amssymb,amsthm}
\usepackage{enumitem}
\usepackage[expansion=false]{microtype}
\usepackage[hidelinks]{hyperref}
\setlist{itemsep=2pt,topsep=4pt}
\newcommand{\E}{\mathbb{E}}
\newcommand{\R}{\mathbb{R}}
\newcommand{\norm}[1]{\left\lVert #1\right\rVert}
\title{Backward Rules and Moving References\\in Continual Test-Time Adaptation\\[0.4em]
\large A research proposal on stability, tracking, and current predictive risk}
\author{Amit Arad\\[0.3em]
\small Advisor: Assaf Shocher (Faculty of Data and Decision Sciences)}
\date{October 2026}

\begin{document}
\maketitle
\begin{abstract}
How does differentiating through a frozen reference's prediction input
interact with a moving reference during test-time adaptation? We propose
to answer this question for a specified smooth surrogate of IT$^3$, first
in a finite-dimensional model and then in one audited implementation.
The provisional target is a condition connecting feedback geometry,
reference update speed, and discrete step size to useful local stability
and tracking. Current predictive risk is the primary objective; retention
and recovery are secondary diagnostics when earlier tasks remain relevant.
Two checked constructions now separate consistency from retention and
demonstrate different local stability for the two backward rules around
correct, nonconstant predictions. The latter has an exact rotational-mode
boundary and a noncommuting example with discrete-time certificates.
These are constructed mechanisms, not benchmark results or novelty
claims. A subsequent five-seed synthetic neural test passed its accuracy
gate but found none of the prescribed strong rotational geometry. A
stationary follow-up now verifies observable finite-time amplification and
risk deterioration under fast EMA. Its predicted uniform ordering of the
two backward rules fails. A local derivative decomposition now shows that
neither residual-weighted terms nor unequal student/reference parameter
Jacobians are required for the measured positive observable modes.
A bounded comparison with the closest teacher-feedback analyses rules
out generic EMA failure as a sufficient paper claim. A paired intervention
now holds the reference's prediction input fixed while preserving its
weight-update rule. It removes at least 88.5\% of excess terminal error
in all four primary seeds, but each cached-input EMA run still deteriorates
relative to its initial accuracy. The next question is whether initial derivatives
can predict the early output separation caused by this input path.
Neither sensitivity nor consistency certifies useful adaptation.
\end{abstract}

\section{Question, setting, and scope}
The value of a consistency loss does not specify its learning dynamics:
freezing a network's parameters and detaching its output can preserve the
same forward values while changing the update direction. A moving reference
adds feedback between the two models. This proposal asks whether that
feedback admits an interpretable stability boundary that predicts useful
adaptation behavior, beyond a generic statement that a linearized system
must be stable.

The learner starts from supervised pretraining and observes an unlabeled
stream $x_t\sim\rho_t$. The adaptive state $\theta_t$ and reference state
$\phi_t$ may persist across batches. Labels are available only to the
evaluator after pretraining. The primary quantity is current task risk
$R_t(\theta)=\E_{(x,y)\sim\mu_t}\ell(q_\theta(x),y)$, evaluated before each
batch's adaptation; post-update risk is reported separately. The initial geometry test uses the source distribution. The first dynamics
test also keeps a fixed empirical distribution stationary; covariate changes with one
unchanged labeling function and compatible recurring domains follow only
if the mechanism is supported. Arbitrary changes in hidden
labels cannot be identified from unlabeled inputs alone.

The local anonymous position manuscript \emph{Forgetting is not Catastrophic,
Our Benchmarks are} motivates evaluating the presently relevant task.
It remains private reading context; its authors and public publication
status are unverified. Its discussion of streams with labels does not
provide an unlabeled-TTA guarantee. Accordingly, old-task retention is an
objective only when those tasks remain relevant. Accuracy loss, recovery
after new adaptation, and preservation of information are different claims.

\section{What the literature and implementation already establish}
IT$^3$ introduces a dual-input predictor and already includes online
adaptation with an EMA reference; its online discussion allows deterioration
on the original distribution~\cite{it3}. Online carry-over and gradient
alignment were also studied in original TTT~\cite{ttt}. Neither online
adaptation nor the observation that aligned gradients improve task loss
is proposed as a new contribution.

The closest continual-adaptation baselines constrain the claim. CoTTA
combines weight/augmentation averaging with stochastic restoration to
pretrained weights~\cite{cotta}; EATA uses sample selection and
importance-weighted regularization against forgetting~\cite{eata}.
RMT studies consistency-loss gradient properties and teacher-based
continual adaptation~\cite{rmt}. PeTTA already considers recurring
environments and analyzes a tractable collapse model~\cite{petta}.
Recent work explicitly compares frozen and moving teachers~\cite{fixedteacher}
and combines dynamic EMA with source anchoring~\cite{desa}. Thus ``EMA can
fail,'' ``a fixed teacher helps,'' and an $A\to B\to A$ experiment are
insufficient paper claims. Nonstationary adaptation guarantees and recovery
complexity also have existing treatments~\cite{realign,recovery}.

\paragraph{What the closest analyses do and do not cover.}
The bounded audit uses PeTTA v4 (Theorem~1 and Appendix~B.3), RMT v2
(Section~3.1, Eqs.~(2)--(6)), and the fixed-teacher preprint v1
(Section~5.2, Eqs.~(3)--(5)). PeTTA bounds a scalar class-mean recurrence;
RMT compares loss derivatives with a fixed teacher probability;
the fixed-teacher paper explicitly describes error feedback through EMA.
These support rejecting generic feedback failure as a distinct contribution,
but do not directly certify our dual-input neural dynamics.
The inspected official RMT and fixed-teacher consistency paths feed images
to the teacher, without a student-prediction input. Exact revisions and code
locations are retained in the research notes; no benchmark was executed.
This limited comparison does not certify novelty~\cite{petta,rmt,fixedteacher}.

\paragraph{Version and graph audit.}
The bibliography uses the ICML 2025 IT$^3$ record. This research pass compared
arXiv v1 and v2; it did not establish byte identity between the supplied
local PDF, arXiv v2, and the proceedings PDF. The official repository was
inspected at commit \texttt{100524001d89}\footnote{Full revision:
\href{https://github.com/nikitadurasov/ittt/tree/100524001d890efc08911f3cba8f45f51f4aa7d1}{\texttt{100524001d890efc08911f3cba8f45f51f4aa7d1}}.}.
The tabular, airfoil, and age regression notebooks pass an undetached student
prediction into a reference excluded from the student optimizer. They use
absolute discrepancies. The CIFAR notebook instead reuses a shared
trainable model across the two adaptation passes and uses Jensen--Shannon
divergence. The linked \texttt{torch-ttt} implementation also reuses its
model. There is therefore no universal ``IT$^3$ backward graph.'' This
proposal selects the regression-style reference-input path and a squared
surrogate explicitly. The inspected repository did not reveal a reproducible
online EMA schedule; our ordering and optimizer-state policy must be declared
as experimental choices. These are static audits, not reproduced results.

\section{The specified dynamics}
Write $q_\theta(x)=f_\theta(x,0)$ and $T_\phi(x,q)=f_\phi(x,q)$. With
$r=T_\phi(x,q_\theta(x))-q_\theta(x)$, define
\begin{equation}
 D(\theta,\phi;\rho)=\tfrac12\E_\rho\norm{r}^2,
 \qquad J_q=\partial_\theta q_\theta(x).
 \label{eq:discrepancy}
\end{equation}
Holding reference parameters fixed during the student derivative retains
the reference-input Jacobian. Detaching the entire reference output removes it:
\begin{align}
 G_{\rm input}&=\E_\rho[J_q^\top(\partial_qT_\phi-I)^\top r],
 \label{eq:input}\\
 G_{\rm stop}&=-\E_\rho[J_q^\top r].
 \label{eq:stop}
\end{align}
These are chain-rule calculations for \eqref{eq:discrepancy}, not results
for every IT$^3$ implementation. $G_{\rm stop}$ is a controlled comparison.
Frozen parameters during differentiation do not imply a reference that
stays fixed over the stream.

We first study $\dot\theta=-G(\theta,\phi;\rho_t)$ and
$\dot\phi=(\theta-\phi)/\tau$, with $\tau>0$. Even for $G_{\rm input}$,
a fixed-distribution calculation gives
\begin{equation}
 \dot D=-\norm{\nabla_\theta D}^2
 +\tau^{-1}\langle\nabla_\phi D,\theta-\phi\rangle.
 \label{eq:energy}
\end{equation}
The second term has no prescribed sign. Distribution drift adds a time
derivative, or a jump in the evaluated objective at a switch. For
$G_{\rm stop}$ the first term is instead
$-\langle\nabla_\theta D,G_{\rm stop}\rangle$.

\paragraph{Discrete algorithm.}
The first study uses gradient descent and a student-first EMA,
\begin{equation}
 \theta_{n+1}=\theta_n-\eta G(\theta_n,\phi_n;\rho_n),\qquad
 \phi_{n+1}=\alpha\phi_n+(1-\alpha)\theta_{n+1},\quad 0<\alpha<1.
 \label{eq:discrete}
\end{equation}
Both states start from the pretrained model. Batch statistics are specified
and held fixed in the first comparison; there is no replay or source-label
access during adaptation. Adam, normalization updates, and alternative EMA
orderings are later robustness checks. Resetting weights without resetting
optimizer state is not treated as an equivalent baseline.

\section{Ranked directions and the provisional leading target}
\begin{enumerate}[leftmargin=*]
\item \textbf{Leading target: coupled dynamics and predictive amplification.}
Determine when reference-input geometry is insufficient to describe the
student--reference parameter dynamics, and whether those dynamics predict
finite-time changes in useful outputs. The constructed stability boundary
is an existence result; its strong rotational mechanism was absent from
the first trained-model screen. Finite-time tests support EMA amplification
but reject a uniform backward-rule ordering. A matrix diagnostic also
rejects the necessity of student/reference derivative mismatch for the
local growing modes. The closest-work audit demotes generic EMA failure
as a paper claim. A realizable cached-input control now removes most
excess terminal error in all four primary seeds, while leaving substantial
residual deterioration. Prioritize an untuned prediction of the resulting
output separation before extending to tracking or claiming a correction. A thesis can develop the audit, qualified theorem, and empirical
test; a first paper requires a distinct mechanism and evidence beyond
existing teacher-adaptation analyses.
\item \textbf{Second: a compatible-task failure and a justified correction.}
Use a diagnosed shared-parameter direction to explain loss on an earlier
still-relevant task at matched current-task progress. Propose a correction
only if its information budget is explicit and it differs substantively
from restoration, anchoring, or protected-gradient methods~\cite{ogd}.
The appendix baseline alone does not meet this threshold.
\item \textbf{Third: recovery and information access.}
Test whether reference lag predicts recovery under a fixed adaptation and
probe budget. This requires separating recovery from relearning and from
already studied representation shifts~\cite{spurious}. It follows only
after the first mechanism is established.
\end{enumerate}

\subsection{A constructed useful equilibrium and a stability distinction}
The earlier absolute model $f_W(x,z)=Wx+A_xz$ has a limitation: strict
stability of its homogeneous dynamics implies zero predictions. To test
whether a useful equilibrium is possible, consider instead the smooth
shared architecture
\begin{equation}
 f_W(x,z)=h(x)+Wx+
 A_x(z-h(x))\frac{h(x)^\top z}{\norm{h(x)}^2},
 \qquad h(x)\ne0.
 \label{eq:realization}
\end{equation}
Here $h$ is a fixed pretrained baseline; $W$ is adapted and $V$ is its EMA.
Then $q_W(x)=f_W(x,0)=h(x)+Wx$, while $f_0(x,h(x))=h(x)$ and
$\partial_z f_0(x,h(x))=A_x$. On two equally likely inputs
$x_1=(1,0)^\top$, $x_2=(3/5,4/5)^\top$, set
$h(x_1)=(1,0)^\top$ and $h(x_2)=(0,1)^\top$, with evaluation labels
equal to $h$. Thus $W=V=0$ represents correct, distinct nonzero
predictions, not collapse. Correctness is planted in the baseline;
the learner receives no labels during adaptation.

Writing $B_x=I-A_x$, the residual is
$Vx-B_xWx+O(\norm{W}^2)$. For column vectorizations $w,v$, the local
Jacobian is
\begin{equation}
 M=\begin{pmatrix}-H&K\\I/\tau&-I/\tau\end{pmatrix},
 \label{eq:block}
\end{equation}
where the two backward rules yield
\begin{align}
 H_{\rm input}&=\E[xx^\top\otimes B_x^\top B_x],&
 K_{\rm input}&=\E[xx^\top\otimes B_x^\top],\nonumber\\
 H_{\rm stop}&=\E[xx^\top\otimes B_x],&
 K_{\rm stop}&=\E[xx^\top\otimes I].
 \label{eq:blocks}
\end{align}
The covariance eigenvalues are $1/5$ and $4/5$, so current risk
$R(W)=\tfrac12\E\norm{Wx}^2$ is positive definite in $W$.
A Hurwitz $M$ therefore gives local exponential recovery to useful
predictions in this construction. Monotone risk is not implied.

\paragraph{An exact boundary and a noncommuting witness.}
Let
\[
 J=\begin{pmatrix}0&-1\\1&0\end{pmatrix},\qquad B_0=\tfrac12 I+J.
\]
For $B_{x_1}=B_{x_2}=B_0$, an exact mode calculation gives local
exponential stability of the input-gradient rule for every finite
$\tau>0$, whereas the stopped-output rule is locally exponentially
stable exactly when $\tau>10$. At $\tau=10$ its linearization has
imaginary roots. This is a linear stability crossing, not a proved
nonlinear Hopf bifurcation. The two rules are consequently more than
a scalar change of speed; noncommutation across inputs is not necessary
for this distinction.

For $B_{x_1}=B_0$ and $B_{x_2}=B_0+\tfrac1{10}\operatorname{diag}(1,-1)$,
the matrices do not commute. Exact rational Routh certificates give the
classifications below. For the student-first discrete update, use
$\eta=1/10$ and $\alpha=\tau/(\tau+\eta)$, with Jacobian
\begin{equation}
 J_{\eta,\alpha}=\begin{pmatrix}
 I-\eta H&\eta K\\
 (1-\alpha)(I-\eta H)&\alpha I+(1-\alpha)\eta K
 \end{pmatrix}.
 \label{eq:discreteblock}
\end{equation}
\begin{center}
\begin{tabular}{lrrr}
Rule & $\tau$ & $\max\Re\lambda(M)$ & $\rho(J_{\eta,\alpha})$\\ \hline
Input gradient & 1 & $-0.09933$ & $0.99024$\\
Stopped output & 1 & $+0.16592$ & $1.01767$\\
Input gradient & 20 & $-0.02134$ & $0.99786$\\
Stopped output & 20 & $-0.02156$ & $0.99782$
\end{tabular}
\end{center}
The displayed eigenvalues are rounded diagnostics. The underlying
continuous classifications use exact characteristic polynomials and
Routh signs; the discrete classifications use the exact bilinear
transformation of the characteristic polynomial and the same test.
The unstable cases have four unstable modes. No interval-wide threshold
is claimed for this noncommuting perturbation. Nonlinear simulations
from matched small student/reference perturbations agree with the
predicted recovery or growth, including a halved-step check.

\paragraph{What the construction establishes.}
A task-correct equilibrium and a directional stability distinction are
mathematically realizable in one smooth dual-input model. The result
does not show that neural pretraining learns this model or that it adapts
successfully to an unseen distribution shift. At exact initialization
it stays fixed; the experiments test recovery from parameter perturbations.
The chosen $A_x$ are not idempotent projections. A genuinely idempotent
differentiable reference on a neighborhood would satisfy $A_x^2=A_x$
at its fixed point, so measured projection-like Jacobians could invalidate
this particular mechanism. Classical feedback stability tests are not
the proposed novelty. Full derivations, exact certificates, and verification
scripts are recorded in the research notes.

\subsection{Assumptions, tracking, and rejection criteria}
A theorem beyond the constructed example must connect a smooth trained
model to a justified task-correct reference point or equilibrium family,
and control the identifiable directions. Any extension to drift must specify how an
equilibrium $z_t^\star$ moves. A bound of the form
$\norm{e(t)}\le C e^{-\gamma t}\norm{e(0)}+
C\sup_s\norm{\dot z_s^\star}/\gamma$
requires a uniform stability estimate and a justified forcing bound;
it is a target template, not a proved result here. Pointwise stable
Jacobians alone do not justify arbitrary switching. Converting state
tracking to predictive risk additionally needs local output regularity
and a task-alignment or correctness assumption.

Reject or narrow the leading direction if the stable state is useless;
if the two rules reduce to time rescaling; if the condition assumes
contraction without explaining its origin; if a correction duplicates
the closest prior work; or if predicted behavior fails inside its stated
regime in small nonlinear tests. A counterexample to a proposed certificate
is a useful result, but should not be relabeled a positive theorem.

\section{Minimal experiment and near-term milestones}
\paragraph{Completed neural geometry gate.}
A fixed five-seed experiment trained a $4\to16\to2$ tanh regressor on
256 source inputs using equal-weight zero/label input supervision inspired
by IT$^3$~\cite{it3}. It used a smooth, nonconstant two-output labeling
function, squared losses, 6000 Adam steps per seed, and 256 independent
evaluation inputs. No architecture search or checkpoint selection was
performed. The fit gate required both auxiliary-input versions to have
normalized mean squared error at most $0.05$, where normalization divides
by total evaluation-label variance.
\begin{center}
\begin{tabular}{rrrr}
Seed & Zero-input NMSE & Label-input NMSE & Fit gate\\ \hline
0 & $0.01203$ & $0.00405$ & pass\\
1 & $0.06620$ & $0.01435$ & fail\\
2 & $0.03434$ & $0.00698$ & pass\\
3 & $0.01708$ & $0.00612$ & pass\\
4 & $0.00719$ & $0.00146$ & pass
\end{tabular}
\end{center}
Four seeds passed, but no evaluation point in any seed met the prescribed
strong-rotation screen, at either predicted or true-label auxiliary input.
The screen required eigenvalues $a\pm i\omega$ with $0<a<1$ and
$\omega^2>a(1-a)$, with a numerical margin; for the normal family this is
the fast-reference separation underlying the construction. For general
nonnormal matrices it is only a screen, not a stability criterion.
The test therefore fails to support that particular explanation in this
model. It does not rule out rotation in other IT$^3$ settings.
Nor does it justify assuming projection geometry: across the four passing
seeds, median relative defects
$\norm{A_x^2-A_x}_F/(\norm{A_x}_F+10^{-12})$ range from $0.46$ to $0.94$.
All sampled prediction-input operator norms are below $0.902$, which does
not imply stability of the coupled parameter updates.

The actual full parameter fields were also differentiated, retaining
nonzero-residual terms. Both rules have positive instantaneous spectral
abscissae at the tested $\tau=1,20$, but the checkpoints have nonzero
velocity: these spectra are not equilibrium stability certificates or
predictions of sustained risk growth. Moreover, 32 of 114 student
parameters multiply the always-zero auxiliary input and receive exactly
zero gradient. With matched initialization those coordinates remain fixed
in both copies; removing them preserves the measured maximal real
eigenvalues. A strict full-parameter Hurwitz target is inappropriate.
Finite-difference gradient checks agree within $2.5\times10^{-10}$, and
halving the full-field differentiation step changes entries by less than
$1.3\times10^{-10}$. Protocol, all checkpoints, derivatives, and results
are retained in the research notes. This is a synthetic mechanism test,
not an official-implementation or distribution-shift reproduction.

\paragraph{Completed finite-time response test.}
The same five checkpoints were tested without retraining; seeds 0,2,3,4
were primary and seed 1 a declared sensitivity case. Adaptation repeatedly
used 64 unlabeled inputs, with a separate 192-input evaluator. This is a
stationary empirical-distribution stress test, not a fresh-sample stream.
Both rules were compared with fixed reference, EMA at $\tau=1,20$, and
no adaptation. The specified student-first update used $\eta=.02$ for
1000 steps ($T=20$), with one halved-step comparison at the same horizon.
Labels entered only the evaluator.

Two student-only initial perturbation directions were fixed per seed:
the student component of a local growing mode, shared across configurations,
and one seeded random direction. Both lie in the 82 non-fixed coordinates;
the reference initially remains unchanged. Prediction sensitivity was
normalized by its initial value for the same direction. The modal probe
deliberately stresses the dynamics; it is not typical noise or a worst-case
singular-value calculation. At $\tau=1$, the primary modal results are:
\begin{center}
\begin{tabular}{rrrrr}
 & \multicolumn{2}{c}{Prediction sensitivity gain}
 & \multicolumn{2}{c}{Final NMSE}\\
Seed & Input gradient & Stopped output & Input gradient & Stopped output\\ \hline
0 & $4.832$ & $5.316$ & $1.703$ & $3.810$\\
2 & $645.165$ & $51.875$ & $54.111$ & $249.731$\\
3 & $23.352$ & $63.814$ & $3.091$ & $9.749$\\
4 & $1.859$ & $1.912$ & $0.516$ & $0.672$
\end{tabular}
\end{center}
Initial primary NMSE ranges from $.00752$ to $.03673$; fixed-reference
final NMSE is $.00685$--$.03697$. The evaluator normalization differs
from the earlier 256-point geometry test. The predefined amplification
gate passes in three of four primary seeds: stopped-output fast-EMA gain
at least 2, with fixed-reference gain at most 1.25. The prediction that
stopped-output gain exceeds input-gradient gain by at least 25\% in three
seeds fails; only seed 3 meets it. Seed 2 illustrates why sensitivity and
risk rankings must remain separate. Slow EMA is less damaging over this
horizon but is not a guarantee: stopped-output seed 2 still deteriorates
materially. Fast-EMA consistency discrepancy increases in these runs;
monotone consistency improvement is not claimed.

The actual time-dependent tangent recurrence predicts paired nonlinear
prediction differences with maximum relative error $1.28\times10^{-5}$
at perturbation size $10^{-5}$. Halving the step changes final modal gains
by at most $2.40\%$ and NMSE by $0.52\%$. These checks support a finite-time
observation, not an equilibrium or ODE-limit theorem. The tangent product
is an established chain-rule tool; its successful implementation is not
novelty. All 30 configurations, including the failed-fit seed, are retained
in the research notes. No distribution drift or retention test was run.

\paragraph{Completed local feedback decomposition.}
At the same five initial checkpoints, write
$J=\partial_\theta f_\theta(x,0)$,
$L=\partial_\phi f_\phi(x,q)$ with $q$ held fixed, and
$B=I-\partial_qT_\phi$. The input-gradient zero-residual blocks of
the student velocity are $-\E[J^\top B^\top BJ]$ and
$\E[J^\top B^\top L]$; the stopped-output blocks are
$-\E[J^\top BJ]$ and $\E[J^\top L]$.
The full derivatives additionally differentiate the prefactor multiplying
the residual. These product-rule remainders were checked independently
against the saved full derivatives, with maximum entry error
$1.58\times10^{-10}$. Dropping them does not make the actual checkpoint
an equilibrium.

The frozen protocol compared full, zero-residual, and additionally
$L=J$ matrices on the 82 non-fixed coordinates per copy, for both rules
and $\tau=1,20$. A positive observable mode requires
$\Re\lambda>10^{-6}$ and
$\norm{J_{\rm eval}u}/(\norm{u}\norm{J_{\rm eval}}_2)\ge10^{-3}$.
Here $u$ is the student component of the mode.
All four primary seeds retain such a mode in every setting under both
diagnostics. Removing residual terms changes the leading observable rate
by at most $7.34\%$, passing the predefined $10\%$ comparison.
Additionally setting $L=J$ changes it by up to $34.12\%$, failing that
quantitative comparison despite preserving positive modes.

Thus neither residual terms nor derivative mismatch are necessary for
this local signature. The $L=J$ replacement is a matrix diagnostic, not
a demonstrated neural intervention; it does not explain the complete
nonlinear risk trajectory or show that mismatch is irrelevant. These
are nonstationary checkpoints, so equilibrium stability is still not
established. All 60 matrices, including the failed-fit seed, and the
independent checks are retained in the research notes.

\paragraph{Completed control: hold the prediction input fixed.}
The reference target $T_\phi(x,q_\theta(x))$ can change with the student
even when $\phi$ is fixed. Detaching its output still recomputes that
target after every update. To remove this forward path, the control
caches $c(x)=q_{\theta_0}(x)$ once and uses
\[
 r_c=T_\phi(x,c(x))-q_\theta(x),\qquad
 \dot\theta=\E[J^\top r_c].
\]
Both backward rules coincide for the control. Its initial target and
velocity equal those of the original stopped-output rule. With
$S=\partial_\theta\dot\theta$, $T=\partial_\phi\dot\theta$ and
$A=\partial_qT_\phi$, the exact initial full-block differences are
$S_{\rm stop}-S_c=\E[J^\top A J]$ and $T_{\rm stop}-T_c=0$.
The residual-weighted prefactor terms cancel at initialization; the
remaining difference has no prescribed sign. Independent differentiation
verified these identities within $6.51\times10^{-11}$.

The unchanged protocol reused all five checkpoints and the 64/192
adaptation/evaluation split. It compared original stopped-output,
original full input-gradient, and cached-input dynamics, each with fixed
reference or $\tau=1$ EMA, at $\eta=.02$ and $.01$, both to $T=20$.
All 60 trajectories completed without retraining or tuning. Labels were
evaluation-only. Every condition allocated the same 128-number prediction
cache (1024 bytes); only the control read it during updates. This is
retained information associated with the repeated inputs, not a proposed
fresh-input online method.

Let $R_0,R_e,R_c$ denote initial, original stopped-output/EMA, and
cached-input/EMA terminal NMSE. The predeclared gate required
$R_e\ge2R_0$, $R_e-R_0\ge.01$, and
$\max(R_c-R_0,0)\le.25(R_e-R_0)$ in at least three of four primary seeds
at both step sizes. All four pass. Main-step results are:
\begin{center}
\begin{tabular}{rrrrr}
Seed & Initial $R_0$ & Original $R_e$ & Cached $R_c$ & Excess removed\\ \hline
0 & $.01234$ & $3.80976$ & $.19615$ & $95.16\%$\\
2 & $.03673$ & $249.73110$ & $.11633$ & $99.97\%$\\
3 & $.01722$ & $9.74948$ & $.07419$ & $99.41\%$\\
4 & $.00752$ & $.67178$ & $.08344$ & $88.57\%$
\end{tabular}
\end{center}
The minimum reduction across both step sizes is $88.55\%$. All original
endpoints exactly reproduce the earlier baseline. Halving the step changes
terminal NMSE by at most $.517\%$ across every condition, below the $5\%$
numerical gate. The failed-fit seed 1 is retained as a sensitivity case
and also meets the numerical/effect inequalities, without entering the
primary count. All states, targets, metrics and derivative checks are
preserved in the research notes.

This supports a large effect of the forward input path under the declared
intervention. It does not establish a safe correction: every cached/EMA
primary run still ends above its initial risk. Cached/EMA discrepancy is
lower at the endpoint despite this accuracy loss; monotonicity is not
claimed. Reference targets remain exactly constant in the fixed-weight,
cached-input control; caching alone leaves weight-driven target changes.
The intervention preserves the reference update rule, not the subsequent
realized reference weights between arms. Neither a general stability
theorem nor novelty follows from the reduction.

\paragraph{Next gate: predict the effect before extending the setting.}
The narrowed thesis question is whether prediction-input feedback admits
an interpretable finite-time prediction separate from reference-weight
tracking and backward detachment. The next saved analysis uses the full
initial derivative blocks and shared nonzero velocity to predict the
early output separation between original and cached-input trajectories.
It fixes $T=2$ and a $25\%$ relative output-error gate in at least three
of four primary seeds at both step sizes, without fitting or rerunning
the nonlinear trajectories. This is a planned analysis of existing data
after terminal results are known, not an independent experiment. A failure
rejects that simple initial-geometry predictor. A pass would still require
a justified approximation bound and robustness to an audited regression
loss and implementation before a first-paper claim. A compatible
$A\to B\to A$ stream remains a later gate.

For that subsequent stream study, compare no adaptation, episodic reset,
carried student with a fixed
reference, carried student with EMA, and the explicit output-detachment
comparison. Select the closest established restoration/anchoring comparator
after the mechanism is identified. Report trainable parameters, reference
copies, source statistics, replay memory, teacher batch statistics, and
update counts; match these budgets or disclose differences. Report current
pre-update risk, post-update risk, discrepancy, reference lag, and boundary
prediction errors. On compatible recurring domains, also report earlier
risk and recovery relative to both no adaptation and a fresh source reset
under equal new-data/update budgets. Compare retention at matched current
risk so that slower learning is not mistaken for protection.

\paragraph{Four milestones, subject to the rejection gates.}
The constructed feedback family, useful fixed point, and discrete
certificates complete the first milestone. The first neural geometry gate
is also complete, with a negative result for strong rotation. The finite-time
gate now supports EMA amplification but rejects uniform backward-rule
ordering. The completed local decomposition rejects mismatch as a
necessary explanation, while the $L=J$ rates are not uniformly
accurate. The closest-prior-art claim map is complete and rules out
generic feedback failure as a sufficient paper claim. The predeclared
prediction-input control is complete and passes in all primary seeds,
while failing to preserve initial accuracy. The next milestone is a
checked prediction of its early output effect; an affine approximation
alone is not a novel theorem. The final milestone is reproduction of one
selected implementation graph and robustness to its loss and optimizer.
At that point, review the closest
papers theorem by theorem and decide whether the result supports a focused
first paper, needs a different mechanism, or remains thesis groundwork.
No publication outcome is assumed.

\paragraph{Existing work and limits.}
The earlier scalar learned-predictor certificate remains a checked local
research artifact, useful for proof practice. Its PEIRA-inspired predictor
tracks an optimizer, whereas the present reference tracks the student;
the earlier result does not prove anything about EMA adaptation. Likewise,
function-space stop-gradient convergence results require their own operator
and regularity assumptions~\cite{stopgrads}. The project currently has a source/code audit, the appendix calculation,
constructed local stability results, and reproducible synthetic neural
geometry, finite-time response, local feedback decomposition, and a
paired frozen-input intervention.
It has no general trained-network stability theorem,
distribution-drift guarantee, official-implementation reproduction,
or verified novelty claim.

\medskip
\noindent{\small\textit{Acknowledgments.} This proposal has developed through ongoing discussion and collaboration with Shoham Grunblat.}

\appendix
\section{A checked affine diagnostic, not a paper claim}
Take unit features $x_A,x_B\in\R^2$ with $0<|c|<1$, where
$c=x_A^\top x_B$. Let $f_\theta(x,z)=\theta^\top x+a(x)z$,
$a(x_A)=0$, and $a(x_B)=1-b$ for $b>1$. Start with
$\theta_0=\phi_0$, let $y_A=\theta_0^\top x_A$ and $y_B=0$, and use
these labels only for evaluation. During the $B$ phase, write
$u=\theta^\top x_B$, $v=\phi^\top x_B$, and $u_0\ne0$. Then
\begin{equation}
 \dot u=k(v-bu),\qquad \tau\dot v=u-v,
 \qquad k=b\ \text{for }G_{\rm input},\quad k=1\ \text{for }G_{\rm stop}.
 \label{eq:toy}
\end{equation}
The two reduced eigenvalues are negative for every $\tau,k>0$; the
orthogonal student component is conserved. In fact $u$ is a positive
weighted sum of two decaying exponentials times $u_0$. Consequently
\begin{equation}
 D=\tfrac12(v-bu)^2\downarrow0,\qquad
 R_B=\tfrac12u^2\downarrow0,\qquad
 R_A=\tfrac12c^2(u_0-u)^2\uparrow\tfrac12c^2u_0^2.
 \label{eq:toyresult}
\end{equation}
Both tasks are jointly realizable. The limiting state is
$\theta_\infty=\phi_\infty=\theta_0-u_0x_B$, independent of the finite
EMA timescale. Returning to $A$ from that limit provides zero consistency
signal and hence no recovery. A finite-time return can recover partially
through reference lag. Full derivations and verification are stored in
the research notes; 40 parameter combinations at 502 times each passed
the closed-form checks, with maximum ODE residual about $1.1\times10^{-13}$.
Numerical checks support the algebra; they are not a neural experiment.

This diagnostic shows that useful current adaptation can coexist with loss
on a compatible earlier task, and that a slow EMA need not preserve it
indefinitely. Its two backward rules only change speed and have identical
retention at equal $B$ progress. The deliberately chosen feedback and label
are not established empirical IT$^3$ properties. The result motivates the
vector study and its falsification gates; it is not an independent novelty
claim. Protecting the old feature direction would require retained
information and overlaps with established projected-update ideas~\cite{ogd}.

\begin{thebibliography}{99}\small
\bibitem{it3} N. Durasov et al. \emph{IT$^3$: Idempotent Test-Time Training}.
ICML, PMLR 267:14867--14883, 2025.
\href{https://proceedings.mlr.press/v267/durasov25a.html}{Published record};
\href{https://arxiv.org/abs/2410.04201}{arXiv:2410.04201, v1/v2}.
\bibitem{ttt} Y. Sun et al. \emph{Test-Time Training with Self-Supervision
for Generalization under Distribution Shifts}. ICML, 2020.
\href{https://proceedings.mlr.press/v119/sun20b.html}{PMLR 119}.
\bibitem{cotta} Q. Wang et al. \emph{Continual Test-Time Domain Adaptation}.
CVPR, 2022.
\href{https://openaccess.thecvf.com/content/CVPR2022/html/Wang_Continual_Test-Time_Domain_Adaptation_CVPR_2022_paper.html}{Primary paper}.
\bibitem{eata} S. Niu et al. \emph{Efficient Test-Time Model Adaptation
without Forgetting}. ICML, 2022.
\href{https://proceedings.mlr.press/v162/niu22a.html}{PMLR 162}.
\bibitem{rmt} M. Doebler, R. A. Marsden, and B. Yang. \emph{Robust Mean
Teacher for Continual and Gradual Test-Time Adaptation}. CVPR, 2023.
\href{https://arxiv.org/abs/2211.13081}{arXiv:2211.13081}.
\bibitem{petta} T.-H. Hoang, D. M. Vo, and M. N. Do. \emph{Persistent
Test-time Adaptation in Recurring Testing Scenarios}. NeurIPS, 2024.
\href{https://arxiv.org/abs/2311.18193}{arXiv:2311.18193}.
\bibitem{fixedteacher} D. Sojka et al. \emph{Rethinking the Teacher-Student
Framework for Test-Time Adaptation}. Preprint, 2026.
\href{https://arxiv.org/abs/2609.02507}{arXiv:2609.02507}.
\bibitem{desa} A. Belal et al. \emph{DESA-TTA: Dynamic EMA and Source
Anchoring for Test-Time Adaptation}. Preprint, 2026.
\href{https://arxiv.org/abs/2609.01795}{arXiv:2609.01795}.
\bibitem{realign} Z.-Y. Zhang et al. \emph{Test-time Adaptation in
Non-stationary Environments via Adaptive Representation Alignment}.
NeurIPS, 2024.
\href{https://papers.neurips.cc/paper_files/paper/2024/hash/abe31a12e83111fdf2cfd54deed5a2ce-Abstract-Conference.html}{Proceedings}.
\bibitem{recovery} Z. Zhou et al. \emph{On the Learnability of Test-Time
Adaptation: A Recovery Complexity Perspective}. ICML, 2026.
\href{https://arxiv.org/abs/2605.28057}{arXiv:2605.28057}.
\bibitem{ogd} M. Farajtabar et al. \emph{Orthogonal Gradient Descent for
Continual Learning}. AISTATS, 2020.
\href{https://proceedings.mlr.press/v108/farajtabar20a.html}{PMLR 108}.
\bibitem{spurious} V. Palit et al. \emph{When Forgetting is not Catastrophic:
On the Mechanics of Spurious Forgetting}. Preprint, 2026.
\href{https://arxiv.org/abs/2610.08718}{arXiv:2610.08718}.
\bibitem{stopgrads} M. W. Shen et al. \emph{How I learned to stop worrying
and love StopGrads: Stationarity, Convergence, and a case study on Flow Map
Learning}. Preprint, 2026.
\href{https://arxiv.org/abs/2609.16222}{arXiv:2609.16222}.
\end{thebibliography}
\end{document}
